% Options for packages loaded elsewhere
\PassOptionsToPackage{unicode}{hyperref}
\PassOptionsToPackage{hyphens}{url}
\documentclass[
  12pt,
]{ctexart}
\usepackage{xcolor}
\usepackage{amsmath,amssymb}
\setcounter{secnumdepth}{-\maxdimen} % remove section numbering
\usepackage{iftex}
\ifPDFTeX
  \usepackage[T1]{fontenc}
  \usepackage[utf8]{inputenc}
  \usepackage{textcomp} % provide euro and other symbols
\else % if luatex or xetex
  \usepackage{unicode-math} % this also loads fontspec
  \defaultfontfeatures{Scale=MatchLowercase}
  \defaultfontfeatures[\rmfamily]{Ligatures=TeX,Scale=1}
\fi
\usepackage{lmodern}
\ifPDFTeX\else
  % xetex/luatex font selection
\fi
% Use upquote if available, for straight quotes in verbatim environments
\IfFileExists{upquote.sty}{\usepackage{upquote}}{}
\IfFileExists{microtype.sty}{% use microtype if available
  \usepackage[]{microtype}
  \UseMicrotypeSet[protrusion]{basicmath} % disable protrusion for tt fonts
}{}
\makeatletter
\@ifundefined{KOMAClassName}{% if non-KOMA class
  \IfFileExists{parskip.sty}{%
    \usepackage{parskip}
  }{% else
    \setlength{\parindent}{0pt}
    \setlength{\parskip}{6pt plus 2pt minus 1pt}}
}{% if KOMA class
  \KOMAoptions{parskip=half}}
\makeatother
\usepackage{longtable,booktabs,array}
\usepackage{calc} % for calculating minipage widths
% Correct order of tables after \paragraph or \subparagraph
\usepackage{etoolbox}
\makeatletter
\patchcmd\longtable{\par}{\if@noskipsec\mbox{}\fi\par}{}{}
\makeatother
% Allow footnotes in longtable head/foot
\IfFileExists{footnotehyper.sty}{\usepackage{footnotehyper}}{\usepackage{footnote}}
\makesavenoteenv{longtable}
\setlength{\emergencystretch}{3em} % prevent overfull lines
\providecommand{\tightlist}{%
  \setlength{\itemsep}{0pt}\setlength{\parskip}{0pt}}
\usepackage{bookmark}
\IfFileExists{xurl.sty}{\usepackage{xurl}}{} % add URL line breaks if available
\urlstyle{same}
\hypersetup{
  hidelinks,
  pdfcreator={LaTeX via pandoc}}

\author{SCX}
\date{2026}

\begin{document}

\section{SCX CodexKnowledge 索引}\label{scx-codexknowledge-ux7d22ux5f15}

\begin{quote}
更新：2026-06-26 \textbar{} SCX v4.0 Phase A
\end{quote}

\subsection{当前状态}\label{ux5f53ux524dux72b6ux6001}

\begin{itemize}
\tightlist
\item
  \textbf{代码}：v4.0 Phase A 完成（通用模块化架构，抽象基类 + YAML
  域配置）
\item
  \textbf{测试}：370 tests 通过
\item
  \textbf{实验}：CIFAR (v1+noise), MedMNIST (v1+v2), AlN v3 (一层+两层)
\item
  \textbf{论文}：5 篇序列（EGP→SCX-Theory→SCX-MLIP→SCX-Sim→SCX-Health）
\item
  \textbf{GitHub}：\texttt{git@github.com:shuchaoxi/State-Conditioned-eXpertise.git}（私有）
\end{itemize}

\begin{center}\rule{0.5\linewidth}{0.5pt}\end{center}

\subsection{文件导航}\label{ux6587ux4ef6ux5bfcux822a}

\subsubsection{核心定义与规划}\label{ux6838ux5fc3ux5b9aux4e49ux4e0eux89c4ux5212}

\begin{longtable}[]{@{}ll@{}}
\toprule\noalign{}
文件 & 用途 \\
\midrule\noalign{}
\endhead
\bottomrule\noalign{}
\endlastfoot
\texttt{SCX\_核心定义.md} & 核心数学定义、符号表、命题索引 \\
\texttt{SCX\_思想扩展\_综合方案.md} & Agent 讨论产出的完整扩展方案 \\
\texttt{SCX\_通用模块化架构设计.md} & v4.0 通用模块化架构设计 \\
\texttt{START\_HERE\_CODEX.md} & 新会话入口（AI 恢复 prompt） \\
\end{longtable}

\subsubsection{TODO（当前唯一有效）}\label{todoux5f53ux524dux552fux4e00ux6709ux6548}

\begin{longtable}[]{@{}ll@{}}
\toprule\noalign{}
文件 & 用途 \\
\midrule\noalign{}
\endhead
\bottomrule\noalign{}
\endlastfoot
\texttt{SCX\_TODO\_2026-06-26.md} & \textbf{当前 TODO} --- Phase 1-4
任务跟踪 \\
\end{longtable}

已归档（\texttt{archive/}）：\texttt{SCX\_TODO.md},
\texttt{SCX\_终极TODO.md}, \texttt{SCX\_TODO\_v4\_通用模块化.md},
\texttt{SCX\_规划\_v4\_通用模块化.md}

\subsubsection{GPT 讨论整理}\label{gpt-ux8ba8ux8bbaux6574ux7406}

\begin{longtable}[]{@{}ll@{}}
\toprule\noalign{}
文件 & 内容 \\
\midrule\noalign{}
\endhead
\bottomrule\noalign{}
\endlastfoot
\texttt{整理\_框架与模块与论文分布.md} & SCX 框架结构 + 5 篇论文分工 \\
\texttt{整理\_IP保护与开源策略与商业模式.md} & IP
保护、开源策略、商业模式 \\
\texttt{整理\_LLM战略与发展路线图.md} & LLM 集成战略 \\
\end{longtable}

\subsubsection{Agent 输出}\label{agent-ux8f93ux51fa}

\begin{longtable}[]{@{}
  >{\raggedright\arraybackslash}p{(\linewidth - 2\tabcolsep) * \real{0.5000}}
  >{\raggedright\arraybackslash}p{(\linewidth - 2\tabcolsep) * \real{0.5000}}@{}}
\toprule\noalign{}
\begin{minipage}[b]{\linewidth}\raggedright
文件
\end{minipage} & \begin{minipage}[b]{\linewidth}\raggedright
内容
\end{minipage} \\
\midrule\noalign{}
\endhead
\bottomrule\noalign{}
\endlastfoot
\texttt{agent\_outputs/01\_SCX\_核心框架\_数学分析.md} & 核心数学框架 \\
\texttt{agent\_outputs/02\_SCX\_子模块\_详细设计.md} & 子模块设计 \\
\texttt{agent\_outputs/03\_竞争分析\_数据估值与主动学习.md} &
竞争分析：数据估值 \\
\texttt{agent\_outputs/04\_竞争分析\_MoE与蒸馏.md} &
竞争分析：MoE/蒸馏 \\
\texttt{agent\_outputs/05\_数学根源与证明.md} & 数学根源 \\
\texttt{agent\_outputs/06\_实现架构与实验设计.md} & 实现架构 \\
\texttt{agent\_outputs/07\_两层描述符\_错误驱动状态发现.md} &
两层描述符方案 \\
\end{longtable}

\subsubsection{原始对话}\label{ux539fux59cbux5bf9ux8bdd}

\begin{longtable}[]{@{}ll@{}}
\toprule\noalign{}
文件 & 说明 \\
\midrule\noalign{}
\endhead
\bottomrule\noalign{}
\endlastfoot
\texttt{与gpt的对话202606261332.md} & 2026-06-26 GPT 对话（9011行） \\
\texttt{与GPT的讨论2026-06-26-07-49.md} & 2026-06-26 早期讨论 \\
\end{longtable}

\subsubsection{其他}\label{ux5176ux4ed6}

\begin{longtable}[]{@{}ll@{}}
\toprule\noalign{}
文件 & 用途 \\
\midrule\noalign{}
\endhead
\bottomrule\noalign{}
\endlastfoot
\texttt{决策日志.md} & 关键决策记录 \\
\texttt{工具状态.md} & 工具配置状态 \\
\end{longtable}

\begin{center}\rule{0.5\linewidth}{0.5pt}\end{center}

\subsection{仓库结构速查}\label{ux4ed3ux5e93ux7ed3ux6784ux901fux67e5}

\begin{verbatim}
SCX/
├── src/scx/                  ← Python 库
│   ├── core/                 在线跟踪、SCX 核心
│   ├── encoders/             ErrorDrivenEncoder, MLIPEncoder, 抽象基类
│   ├── state/                两层状态发现、鲁棒性
│   ├── valuation/            数据分类、自适应阈值、影响力
│   └── action/               压缩、路由
├── tests/                    370 tests
├── theory/
│   ├── definitions/          核心数学定义 + 符号表
│   └── propositions/         Prop 01-06（含两层描述符）
├── experiments/
│   ├── cifar/                CIFAR-10/100 实验
│   ├── mlip_case/            AlN v3 SCX 分析
│   ├── ml_benchmarks/        通用 ML 基准
│   └── synthetic/            合成数据
├── scx-health/               MedMNIST 实验
│   ├── experiments/          compress/noise/routing 脚本
│   └── results/              v1 + v2 报告
├── paper/                    论文草稿
└── CodexKnowledge/           ← 本目录（知识管理）
\end{verbatim}

\subsection{当前优先级}\label{ux5f53ux524dux4f18ux5148ux7ea7}

\begin{enumerate}
\def\labelenumi{\arabic{enumi}.}
\tightlist
\item
  \textbf{P0}：Compress 实验修复（GPU + ResNet-18，恢复 v1 的 clean
  win）
\item
  \textbf{P1}：Noise 实验用状态相关噪声（非均匀）
\item
  \textbf{P2}：Benchmark suite + Online SCX
\item
  \textbf{论文}：Paper 4 (SCX) §4.3 压缩定理 + §6.2 通用实验
\end{enumerate}

\end{document}
